\section{Related Work}
\label{sec:related}

\textbf{Object-level and scene-graph maps.}
Fusion++ maintains object-level TSDF volumes with an existence probability per object~\cite{fusionpp2018}.
3D scene graphs organise objects, places and rooms in layers~\cite{dsg3d2020,hydra2022}, and open-vocabulary systems such as ConceptGraphs associate 2D masks into 3D objects by thresholded geometric and semantic similarity~\cite{conceptgraphs2024}.
These systems mainly target static scenes; apart from existence probabilities as in Fusion++, a stale object is usually fused with new evidence or kept.

\textbf{Change-aware and long-term mapping.}
Persistence filters model the survival time of features in semi-static environments~\cite{rosen2016persistence}, and Perpetua extends this to emergence and persistence of objects~\cite{perpetua2025}.
Volumetric approaches detect object-level changes and update submaps~\cite{pocd2022,panopticmultitsdf2022}, and Khronos estimates the spatio-temporal history of a scene~\cite{khronos2024}.
Closest to our setting, DSG builds dynamic scene graphs in changing indoor environments with render-and-compare removal and introduces Dyn-THOR, with node precision, recall and F1 at 3D IoU 0.3 and the Missing residual rate~\cite{dsg2026}; DynamicGSG updates a Gaussian scene graph by rendering~\cite{dynamicgsg2025}.
DSG and DynamicGSG update with hand-set rules, and their removal takes the object out of the map; persistence models~\cite{rosen2016persistence,perpetua2025} and spatio-temporal maps~\cite{khronos2024} keep information about objects that disappear and re-emerge, through probabilistic models or optimisation rather than costs learned from labels.
We keep the metric definitions of Dyn-THOR where possible, implement their core mechanisms as adapters on our shared front end (Section~\ref{sec:arms}), and ask a different question: whether learned, reversible lifecycle operations add value over learned association alone.
Our adapters are independent implementations of published mechanisms, not reproductions of these systems.

\textbf{Learned memory operations.}
LLM memory systems choose operations such as add, update and delete with a frozen language model~\cite{mem0_2025}; we report a zero-training LLM operator on the same sealed feature tables only in the appendix.
Training on the model's own memory trajectories follows the data-aggregation idea of DAgger~\cite{dagger2011}.
As in work that pairs frozen pre-trained features with a small learned module~\cite{dinowm2025}, the front end (DINOv2 descriptors~\cite{dinov2_2024}, depth geometry and masks) is frozen, so differences between arms come from the memory logic alone.
